\section{Teorema di Cook-Levin}
\[\CNFSAT \text{ é } \NP\text{-completo}.\]

\subsection{$\CNFSAT \in \NP$}
\begin{proof}
Ricordiamo che abbiamo descritto un possibile algoritmo non deterministico
per $\SAT$ giá in \ref{problema:sat}. L'idea era di usare un semplice algoritmo
di model checking su un modello scelto non deterministicamente. \\

In sostanza: \begin{enumerate}
\item Si genera non deterministicamente un modello $M$ con \textbf{choose}.
\item Si verifica se $M \models x$ con un algoritmo di model checking in tempo polinomiale.\\
\end{enumerate}

Descriviamo piú formalmente la MdT, sfruttando il fatto che $x$ sia in CNF. \\
In una CNF ogni clausola deve essere vera, e una clausola é vera se almeno uno dei suoi letterali é vero.
\begin{enumerate}
\item Scansione di $x$ per trovare tutte le proposizioni atomiche $p_1, \ldots, p_n \in AP(x)$.
\item Per ogni proposizione atomica $p_i$, si genera non deterministicamente un suo valore di verità. 
Se la scelta non deterministica porta a $p_i$ vero, allora si scrive $p_i$ sul nastro, altrimenti non si scrive nulla.
\item Si scandisce di nuovo $x$, scansionando $M$ a ogni passo per verificare se soddisfa $p_i$, il letterale sotto la testina. \\
    Essendo $x$ in CNF, se $p_i$ é vero si passa direttamente alla prossima clausola, altrimenti si passa al prossimo letterale.
\item Se si arriva alla fine di una clausola senza trovare un letterale vero, allora la clausola é falsa e si rifiuta.
\item Se si arriva alla fine di tutta la formula (quindi tutte le clausole erano vere), si accetta.
\end{enumerate}


Con il non determinismo, la prima parte ha costo pari a $|M|$, 
perché non bisogna fare altro che scrivere il modello sul nastro.
Il modello viene rappresentato scrivendo solo le proposizioni atomiche che sono vere in $M$, 
e $M \subseteq AP(x)$ implica che $|M| \leq |x|$, quindi il costo di spazio e tempo della prima parte é lineare in $|x|$. \\

La seconda parte ha costo quadratico: scansione di due nastri di dimensione $|x|$.\\

Quindi l'intero algoritmo é polinomiale.
\end{proof}


\newpage
\subsection{$\CNFSAT$ é $\NP$-hard}
$\forall L \in \NP$, $L$ puó essere ridotto a $\CNFSAT$.
\begin{proof}
Sia $N$ una qualsiasi MdT non deterministica a un nastro, che accetta $x \in L$ e opera in tempo polinomiale $f(n)$.\\
Dobbiamo programmare una MdT deterministica che opera in tempo polinomiale e spazio logaritmico,
sullo stesso input $x$, e che riduce la computazione di $N$ a una formula in CNF, $R(x)$,
tale che $R(x)$ é soddisfacibile se e solo se $N$ accetta $x$.\\

Usiamo ancora un altro tipo di rappresentazione della computazione di $N$: 
dopo un albero nel teorema \ref{theorem:non_deterministic_to_deterministic} 
e un grafo nel teorema \ref{theorem:non_deterministic_space_to_deterministic_time},
ora usiamo una tabella, detta \textbf{computation table}.\\

La computation table é una tabella quadrata di lato $T = f(n)$.\\

La riga $t$ rappresenta lo stato della MdT al passo $t$ della computazione,
e la colonna $s$ rappresenta il contenuto della cella $s$ del nastro.\\

Poiché lo spazio é limitato superiormente dal tempo, i caratteri su una riga non
saranno piú di $T$. Se sono di meno, riempiamo la riga con ${\sqcup}$. \\

Rappresentiamo contemporaneamente lo stato e la posizione della testina,
segnando lo stato a pedice della cella in cui si trova la testina.\\

\[
\begin{array}{lllllll}
{\triangleright}_s  & 0    & 1          & 0          & {\sqcup}    & {\sqcup} & {\sqcup} \\
{\triangleright}    & 0_s  & 1          & 0          & {\sqcup}    & {\sqcup} & {\sqcup} \\
{\triangleright}    & 0    & 1_s        & 0          & {\sqcup}    & {\sqcup} & {\sqcup} \\
{\triangleright}    & 0    & 1          & 0_s        & {\sqcup}    & {\sqcup} & {\sqcup} \\
{\triangleright}    & 0    & 1          & 0          & {\sqcup}_s  & {\sqcup} & {\sqcup} \\
{\triangleright}    & 0    & 1          & 0_q        & 1           & {\sqcup} & {\sqcup} \\
{\triangleright}    & 0    & 1          & 1_{"yes"}  & 1           & {\sqcup} & {\sqcup}
\end{array}
\]
\text{ }\\
$R(X)$ descriverá la computation table di $N$.\\

Definiamo le seguenti variabili proposizionali, dove $s$ indica la riga e $t$ la colonna della computation table
($s, t$ vanno da $1$ a $T$):
\begin{itemize}
    \item $Q_t^q \iff $ al passo $t$ lo stato é $q \in Q$
    \item $H_{s,t} \iff $ la testina é sulla cella $s$ al passo $t$
    \item $C_{s,t}^\sigma \iff $ la cella $s$ al passo $t$ contiene il simbolo $\sigma \in \Gamma$
\end{itemize}
\newpage
$R(x)$ sará la congiunzione delle seguenti formule:
\begin{itemize}
    \item Ogni passo puó avere uno e un solo stato
    \[
    \bigwedge_{t=1}^{T} \left( \left( \bigvee_{q \in Q} Q_t^q \right) \land \left( \bigwedge_{q \in Q} \bigwedge_{q' \neq q} (\neg Q_t^q \lor \neg Q_t^{q'}) \right) \right)
    \]
    $T$ clausole sulla presenza di uno stato ($\bigvee$ forma un'unica clausola) e $T \cdot |Q|^2$ clausole per la sua unicitá.\\
    
    \item Ad ogni passo, la testina é su una e una sola cella
    \[
    \bigwedge_{t=1}^{T} \left( \left(\bigvee_{s=1}^{T} H_{s,t}\right) \land \left( \bigwedge_{s=1}^{T} \bigwedge_{s' \neq s} (\neg H_{s,t} \lor \neg H_{s', t}) \right) \right)
    \]
    $T$ clausole sulla posizione della testina e $T^3$ clausole per la sua unicitá in ogni passo.\\
    
    \item Ad ogni passo, ogni cella contiene uno e un solo simbolo
    \[
    \bigwedge_{t=1}^{T} \bigwedge_{s=1}^{T} \left( \left( \bigvee_{\sigma \in \Gamma} C_{s,t}^\sigma \right) \land \left( \bigwedge_{\sigma \in \Gamma} \bigwedge_{\sigma' \neq \sigma} (\neg C_{s,t}^\sigma \lor \neg C_{s,t}^{\sigma'}) \right) \right)
    \]
    $T^2$ clausole per la presenza di un simbolo in ogni cella e $T^2 \cdot |\Gamma|^2$ clausole per l'unicitá del simbolo in ogni cella.\\
    
    \item Ad ogni passo, la cella 1 contiene sempre il simbolo $\triangleright$
    \[
    \bigwedge_{t=1}^{T} C_{1,t}^{\triangleright}
    \]
    $T$ clausole.\\
    
    \item Per ogni passo $t$ tranne l'ultimo, l'unica cella che puó cambiare nel passo successivo $t+1$ é quella su cui si trova la testina al passo $t$.\\
    Cioé, se la testina non si trova in $s$, e su $s$ c'é scritto $\sigma$, allora al prossimo passo ci sará ancora scritto $\sigma$.
    $(\neg H_{s,t} \land C_{s,t}^\sigma) \rightarrow C_{s,t+1}^\sigma$. 
    Per trasformarla in CNF, neghiamo l'antecedente, usando De Morgan, e trasformiamo l'implicazione in OR:
    \[
    \bigwedge_{t=1}^{T-1} \bigwedge_{s=1}^{T} \bigwedge_{\sigma \in \Gamma} \left( H_{s,t} \lor \neg C_{s,t}^\sigma \lor C_{s,t+1}^\sigma \right)
    \]
    É abbastanza intuitivo anche cosí: o la testina si trova su $s$ (e quindi puó cambiare), o su $s$ non é scritto $\sigma$, oppure ci sará scritto $\sigma$ anche dopo.\\
    $T^2 \cdot |\Gamma|$ clausole.\\

    \item Al primo passo, la macchina inizia allo stato $s$, con la testina sulla cella $1$, e il nastro contiene l'input $x = (\sigma_2, \dots, \sigma_{n+1})$ seguito da ${\sqcup}$.
    \[
    Q_1^s \land H_{1,1} \land \bigwedge_{s=2}^{n+1} C_{s, 1}^{\sigma_{s}} \land \bigwedge_{s=n+2}^{T} C_{s, 1}^{\sqcup}
    \]
    Questa formula é particolarmente importante, perché ci dice qual é l'input.\\
    $2 + T$ clausole.\\
    
    \item Le configurazioni \textbf{non} cambiano se non come specificato dalla relazione di transizione $\Delta$.
    Intuitivamente, per ogni quintupla $(q, \sigma, q', \sigma', dir) \notin \Delta$, fissati $s, t$, sia 
    $s' = \begin{cases}
        s-1 & \text{se } dir = \leftarrow \\
        s+1 & \text{se } dir = \rightarrow\\
        s & \text{se } dir = -
    \end{cases}$\\
    $Q_t^q \land H_{s,t} \land C_{s,t}^\sigma \rightarrow \neg (Q_{t+1}^{q'} \land H_{s',t+1} \land  C_{s,t+1}^{\sigma'})$
    che trasformiamo in CNF come prima:
    \[
    \bigwedge_{t=1}^{T-1} \bigwedge_{s=1}^{T} \bigwedge_{(q, \sigma, q', \sigma', dir) \notin \Delta} \left( \neg Q_t^q \lor \neg H_{s,t} \lor \neg C_{s,t}^\sigma \lor \neg Q_{t+1}^{q'} \lor \neg C_{s,t+1}^{\sigma'} \lor \neg H_{s',t+1} \right)
    \]
    Questa formula descrive la computazione di $N$, quindi é cruciale, insieme a quella precedente, per la correttezza della riduzione.\\
    Attenzione solo al fatto che abbiamo "barato" con $s'$, in realtá ci vorrebbero tre formule di questo tipo,
    considerando tutte e tre le possibili direzioni e sostituendo ogni volta $s'$ con la cella indicata.
    Ma asintoticamente non cambia nulla.\\
    $O(T^2 \cdot |Q|^2 \cdot |\Gamma|^2 \cdot 9)$ clausole.\\

    \item Al passo $T$, la macchina deve essere in uno stato di accettazione. 
    Senza questa formula, sarebbero soddisfacibili anche computazioni che terminano in $"no"$.\\
    \[
    Q_T^{"yes"}
    \]
    Il conteggio delle clausole in questa formula é lasciato al lettore per esercizio.\\
\end{itemize}

In totale, ricordando che $|Q|$ e $|\Gamma|$ sono costanti,
il numero di clausole é polinomiale in $T=f(n)$, 
dove $f(n)$ per ipotesi é polinomiale su $n$, 
e quindi lo é anche $|R(x)|$ 
per la chiusura dei polinomi rispetto alla composizione.\\

\textbf{Correttezza}\\
Dobbiamo dimostrare che $N$ accetta $x$ se e solo se $R(x)$ é soddisfacibile.\\
\begin{itemize}
    \item $\Rightarrow$ Se $N$ accetta $x$, allora costruiamo la formula $R(x)$ che riduce la sua computation table.
    Le condizioni sono tutte soddisfatte: la MdT si comporta come dovrebbe, 
    le configurazioni seguono la sua relazione di transizione, l'input é corretto
    e la computazione accetta l'input.\\
    Viceversa, se $N$ non accetta $x$, la condizione finale non puó essere soddisfatta, quindi $R(x)$ non é soddisfacibile.\\
    \item $\Leftarrow$ Se $R(x)$ é soddisfacibile, allora esiste un modello $M$ che soddisfa tutte le condizioni.\\
    In particolare, il modello soddisfa il comportamento corretto della macchina,
    la sua relazione di transizione, l'input corretto e il fatto che la computazione viene accettata.\\
    Viceversa, se non é soddisfacibile, deve esserci una contraddizione in una di queste condizioni,
    ma questo vorrebbe dire che qualsiasi macchina che si comporta come dovrebbe, dato l'input e la relazione di transizione,
    non accetta l'input, e che quelle che accettano alla fine o hanno un diverso input, o una diversa relazione di transizione,
    oppure sono mal formate.\\
\end{itemize}

\textbf{Efficienza}\\
La formula é generata come nella riduzione di Hamilton Path a $\CNFSAT$ (\ref{riduzione:hamiltonpath_sat}).
Si itera su tutte le variabili in gioco: gli stati e l'alfabeto di $N$, 
le celle, i passi e le posizioni della testina nella sua computazione.\\
Il numero di iterazioni é polinomiale.\\

Si usano contatori per generare una clausola alla volta, 
copiandola sul nastro di output se vogliamo che sia inclusa nella formula $R(x)$, 
e poi sovrascrivendola con la successiva.\\
La lunghezza dei contatori é logaritmica rispetto al numero di iterazioni,
e per un polinomio di grado $d$ vale che $O(\log(n^d)) = O(d \log(n)) = O(\log(n))$
quindi lo spazio é logaritmico rispetto alla lunghezza dell'input.\\
\end{proof}